\documentclass[11pt,a4paper]{article}
\usepackage[utf8]{inputenc}
\usepackage[T1]{fontenc}
\usepackage{fouriernc}
\usepackage{tikz}
\usetikzlibrary{arrows.meta}
\usepackage[french]{babel}
\frenchbsetup{StandardLists=true}
\NoAutoSpaceBeforeFDP
\usepackage{enumitem}
\usepackage{amsmath}
\usepackage{amssymb}
\usepackage[left=2cm,right=2cm,top=2cm,bottom=2cm]{geometry}
\usepackage[dvipsnames]{xcolor}
\usepackage{graphicx}
\usepackage{setspace}
\singlespacing
\usepackage{tcolorbox}
\usepackage{multicol}
\usepackage{multirow}
\usepackage{tabularx}
\usepackage{booktabs}
\usepackage{pifont}
\usepackage{array}
\usepackage{fancyhdr}
\setlength{\headheight}{14pt}

\pagestyle{fancy}
\renewcommand\headrulewidth{1pt}
\fancyhead[L]{Lycée Denis Diderot}
\fancyhead[R]{Classe : T~Micro}
\setcounter{page}{1}
\fancyfoot[C]{Les moteurs --- Corrigé --- page~\thepage}
\fancyfoot[R]{\today}
\fancyfoot[L]{Alexis RUIZ}

\setlength{\parindent}{0pt}
\renewcommand{\arraystretch}{1.8}

\newcommand*\acocher{\leavevmode{\fboxsep0pt \fboxrule0.8pt
  \fbox{\vrule height1.8ex depth.1ex width0pt \vrule height0pt depth0pt width1.9ex }}}

\newcommand*\bonne{\leavevmode{%
  \fboxsep0pt \fboxrule0.8pt
  \fcolorbox{green!60!black}{green!20}{%
    \vrule height1.8ex depth.1ex width0pt
    \hspace{0.5pt}\textcolor{green!50!black}{\bfseries\ding{51}}\hspace{0.5pt}}}}

\newcommand{\reponse}[1]{%
  \begin{tcolorbox}[colback=green!5!white, colframe=green!50!black,
                    left=2mm, right=2mm, top=1mm, bottom=1mm]
  #1
  \end{tcolorbox}}

\begin{document}

\begin{center}
  \LARGE{\textbf{Les moteurs --- Corrigé}}
  \vspace{3mm}
\end{center}

\hrule
\vspace{4mm}

%%%%%%%%%%%%%%%%%%%%%%%%%%%%%%%%%%%%%%%%%%
\textbf{Exercice 1 : QCM} \hfill \textbf{/4 points}

\vspace{2mm}

La bonne réponse est indiquée par \bonne{}.

\vspace{3mm}

\begin{center}
\renewcommand{\arraystretch}{2.5}
\begin{tabular}{|m{4mm}|p{4.5cm}||m{3.5cm}|m{3.5cm}|m{3.5cm}|}
	\hline
	\ding{202} & Le moteur électrique est un convertisseur qui transforme :
		& \acocher~de l'énergie mécanique en énergie électrique
		& \bonne~de l'énergie électrique en énergie mécanique
		& \acocher~de l'énergie thermique en énergie électrique \\
	\hline
	\ding{203} & Le rendement d'un moteur est :
		& \acocher~toujours supérieur à 1
		& \acocher~toujours égal à 1
		& \bonne~toujours inférieur à 1 \\
	\hline
	\ding{204} & Pour faire varier la fréquence de rotation d'un moteur asynchrone, il faut faire varier :
		& \acocher~la résistance de l'induit
		& \bonne~la fréquence de la tension d'alimentation
		& \acocher~uniquement l'intensité du courant \\
	\hline
	\ding{205} & La fréquence de rotation d'un moteur peut s'exprimer en :
		& \acocher~Hertz (Hz)
		& \bonne~tours par minute (tr/min) ou rad/s
		& \acocher~Watts (W) \\
	\hline
	\ding{206} & La tension d'alimentation d'un moteur s'exprime en :
		& \acocher~Ampères (A)
		& \acocher~Ohms ($\Omega$)
		& \bonne~Volts (V) \\
	\hline
	\ding{207} & Sur la plaque signalétique d'un moteur, on peut lire :
		& \acocher~le prix d'achat du moteur
		& \bonne~la puissance utile, la tension et la vitesse de rotation
		& \acocher~uniquement le type de moteur \\
	\hline
	\ding{208} & La force électromotrice $E$ d'un moteur à courant continu est :
		& \acocher~indépendante de la vitesse de rotation
		& \acocher~toujours égale à la tension d'alimentation $U$
		& \bonne~proportionnelle à la vitesse de rotation $n$ \\
	\hline
\end{tabular}
\end{center}

\vfill
\newpage

%%%%%%%%%%%%%%%%%%%%%%%%%%%%%%%%%%%%%%%%%%
\textbf{Exercice 2 : Comment vérifier la compatibilité d'un moteur ?}
\hfill \textbf{/7 points}

\vspace{3mm}

\textbf{1.} \hfill \textit{(0,5~pt)}

\reponse{Il s'agit d'un moteur \textbf{asynchrone triphasé} (ASYNCHRONE / 3$\sim$ / triangle).}

\textbf{2.} \hfill \textit{(1,5~pt)}

\reponse{
\renewcommand{\arraystretch}{1.4}
\begin{tabular}{lll}
	\textbf{Grandeur} & \textbf{Symbole} & \textbf{Unité} \\
	Tension d'alimentation  & $U$              & Volt (V) \\
	Intensité du courant    & $I$              & Ampère (A) \\
	Facteur de puissance    & $\cos\varphi$    & sans unité \\
	Puissance utile         & $P_{\rm utile}$  & Watt (W) ou kilowatt (kW) \\
	Vitesse de rotation     & $n$              & tours par minute (tr/min) \\
\end{tabular}
}

\textbf{3.} \hfill \textit{(0,5~pt)}

\reponse{Le technicien doit choisir un moteur du \textbf{même type} (asynchrone), avec une \textbf{puissance utile au moins égale} ($P_{\rm utile} \geq 4{,}2$~kW) et la \textbf{même vitesse de rotation} ($n = 1350$~tr/min).}

\textbf{4.} \hfill \textit{(1,5~pt)}

\reponse{
\[P_a = \sqrt{3} \times U \times I \times \cos\varphi = \sqrt{3} \times 400 \times 10{,}7 \times 0{,}78\]
\[\boxed{P_a \approx 5~780~\text{W} = 5{,}78~\text{kW}}\]
}

\textbf{5.} \hfill \textit{(1,5~pt)}

\reponse{
\[\eta = \frac{P_{\rm utile}}{P_a} = \frac{4~200}{5~780}\]
\[\boxed{\eta \approx 0{,}727 = 72{,}7~\%}\]
}

\textbf{6.} \hfill \textit{(1~pt)}

\reponse{
Le moteur à remplacer : $P_{\rm utile} = 4{,}2$~kW et $\eta \approx 73~\%$.\\[2mm]
\textbf{Moteur 1} : type courant continu $\rightarrow$ \textbf{incompatible}.\\
\textbf{Moteurs 2 et 3} : puissance insuffisante (4,0 et 4,1~kW $<$ 4,2~kW) $\rightarrow$ \textbf{non conformes}.\\
\textbf{Moteur 4} : asynchrone, $P_{\rm utile} = 4{,}2$~kW, $\eta = 78~\%$, $n = 1350$~tr/min $\rightarrow$ \textbf{le bon choix}.
}

\textbf{7.} \hfill \textit{(0,5~pt)}

\reponse{Pour choisir un moteur de remplacement, il faut vérifier que son \textbf{type}, sa \textbf{puissance utile} et sa \textbf{vitesse de rotation} sont compatibles avec le moteur d'origine. On privilégiera un moteur avec un \textbf{meilleur rendement} pour réduire la consommation électrique.}

\vfill
\newpage

%%%%%%%%%%%%%%%%%%%%%%%%%%%%%%%%%%%%%%%%%%
\textbf{Exercice 3 : Rendement d'un moteur asynchrone}
\hfill \textbf{/4 points}

\vspace{4mm}

\textbf{1.} Calcul de la puissance électrique absorbée $P_e$ \hfill \textit{(2~pts)}

\reponse{
Puissance transmise au rotor :
\[P_{\rm tr} = P_m + \text{pertes rotor} = 14~500 + 855 = 15~355~\text{W}\]
Puissance électrique absorbée :
\[P_e = P_{\rm tr} + \text{pertes stator} = 15~355 + 1~745\]
\[\boxed{P_e = 17~100~\text{W} = 17{,}1~\text{kW}}\]
}

\textbf{2.} Calcul des pertes mécaniques \hfill \textit{(2~pts)}

\reponse{
La puissance utile :
\[P_u = \eta \times P_e = 0{,}81 \times 17~100 = 13~851~\text{W}\]
Les pertes mécaniques :
\[P_{\rm mec} = P_m - P_u = 14~500 - 13~851\]
\[\boxed{P_{\rm mec} \approx 649~\text{W}}\]
}

\vfill
\newpage

%%%%%%%%%%%%%%%%%%%%%%%%%%%%%%%%%%%%%%%%%%
\textbf{Exercice 4 : Fréquence de rotation et force électromotrice}
\hfill \textbf{/5 points}

\vspace{4mm}

\textbf{1.} \hfill \textit{(1~pt)}

\reponse{
\[E = U - r \cdot I = 100 - 0{,}9 \times 12{,}1 = 100 - 10{,}89\]
\[\boxed{E \approx 89~\text{V}}\]
}

\textbf{2.} Tableau complété \hfill \textit{(2~pts)}

\vspace{3mm}

\begin{center}
\renewcommand{\arraystretch}{2.2}
\begin{tabular}{|l|m{2.2cm}|m{2.2cm}|m{2.2cm}|m{2.2cm}|m{2.2cm}|}
	\hline
	$U$ (en~V) & \centering 100 & \centering 161 & \centering 242 & \centering 324 & \centering\arraybackslash 407 \\
	\hline
	$I$ (en~A) & \centering 12,1 & \centering\textcolor{green!50!black}{\textbf{23,3}} & \centering 34 & \centering 44,1 & \centering\arraybackslash 51,8 \\
	\hline
	$n$ (en~tr/min) & \centering 500 & \centering\textcolor{green!50!black}{\textbf{800}} & \centering\textcolor{green!50!black}{\textbf{1~200}} & \centering\textcolor{green!50!black}{\textbf{1~600}} & \centering\arraybackslash\textcolor{green!50!black}{\textbf{2~000}} \\
	\hline
	$E$ (en~V) & \centering\textcolor{green!50!black}{\textbf{89}} & \centering 140 & \centering\textcolor{green!50!black}{\textbf{211}} & \centering\textcolor{green!50!black}{\textbf{284}} & \centering\arraybackslash 361 \\
	\hline
\end{tabular}
\end{center}

\vspace{2mm}
{\small\textit{$E = U - 0{,}9 \times I$ \quad;\quad $n$ déduit de la proportionnalité $E \propto n$}}

\vspace{4mm}

\textbf{3.} Graphique $E = f(n)$ \hfill \textit{(1,5~pt)}

\vspace{3mm}

\begin{center}
\begin{tikzpicture}[scale=1]
	\draw[very thin, gray!30] (0,0) grid[xstep=0.2, ystep=0.2] (10,7);
	\draw[thin, gray!60] (0,0) grid[xstep=2, ystep=1] (10,7);
	\draw[->, thick] (0,0) -- (10.4,0) node[right] {$n$ (tr/min)};
	\draw[->, thick] (0,0) -- (0,7.4) node[above] {$E$ (en~V)};
	\foreach \x/\lbl in {2/500, 4/1000, 6/1500, 8/2000, 10/2500}{
		\draw[thick] (\x, 0.08) -- (\x, -0.08);
		\node[below, font=\small] at (\x, 0) {\lbl};
	}
	\foreach \y/\lbl in {1/50, 2/100, 3/150, 4/200, 5/250, 6/300, 7/350}{
		\draw[thick] (0.08, \y) -- (-0.08, \y);
		\node[left, font=\small] at (0, \y) {\lbl};
	}
	\node[below left, font=\small] at (0,0) {0};
	% Droite E = 0,178 n  (n en cm : 2cm=500tr/min ; E en cm : 1cm=50V)
	\draw[red!70!black, very thick] (0,0) -- (7.85, 7.0);
	% Points : (500,89)->(2,1.78) ; (800,140)->(3.2,2.80) ; (1200,211)->(4.8,4.22) ; (1600,284)->(6.4,5.68) ; (2000,361)->(8.0,7.22->clamp 7.0)
	\foreach \pt in {(2,1.78),(3.2,2.80),(4.8,4.22),(6.4,5.68),(8.0,7.0)}{
		\fill[blue!70!black] \pt circle (3pt);
	}
\end{tikzpicture}
\end{center}

\textbf{4.} \hfill \textit{(0,5~pt)}

\reponse{Les 5 points sont \textbf{alignés sur une droite passant par l'origine}. On conclut que la force électromotrice $E$ est \textbf{proportionnelle} à la fréquence de rotation $n$ : $\boldsymbol{E = k \cdot n}$ où $k$ est une constante caractéristique du moteur.}

\vfill

\begin{center}
  \large{\textbf{--- Barème total : /20 points ---}}
\end{center}

\end{document}
